\begin{abstract}    
    \Glspl{mdlm} are trained to in-fill positions in randomly masked sequences, in contrast to traditional \gls{ntp} models. Discussions around \glspl{mdlm} focus on two benefits: (1) any-order decoding  and 2) multi-token decoding. However, we observe that for math and coding tasks, any-order algorithms often underperform or behave similarly to \textit{left-to-right sampling}, and standard multi-token decoding significantly degrades performance. At inference time, \glspl{mdlm} compute the conditional distribution of all masked positions. A natural question is: \textit{How can we justify this additional compute when left-to-right one-token-at-a-time decoding is on par with any-order decoding algorithms?} These findings warrant rethinking how \glspl{mdlm} are utilized. 
    First, we propose \textit{reasoning-as-infilling}. By using \glspl{mdlm} to infill a reasoning template, we can structure outputs and distinguish between reasoning and answer tokens. In turn, this enables measuring answer uncertainty \textit{during} reasoning, and early exits when the model converges on an answer. Next, given an answer, \textit{reasoning-as-infilling} enables sampling from the \gls{mdlm} posterior over reasoning traces \textit{conditioned on the answer}, providing a new source of high-quality data for post-training.     
    On \textsc{gsm8}k, we observe that fine-tuning LLaDA-8B Base on its posterior reasoning traces provides a performance boost on par with fine-tuning on human-written reasoning traces.     
    Additionally, given an answer, reasoning-as-infilling provides a method for scoring the correctness of the reasoning process at intermediate steps, without requiring expensive rollouts or an external model.  
    Second, we propose \gls{med}, a simple adaptive sampler that minimizes the error incurred by decoding positions in parallel based on the conditional entropies of those positions. \gls{med} preserves performance across benchmarks and leads to $2.7\times$ fewer steps. Combined with early exits, \gls{med} leads to a $3.3\times$ speed-up on \textsc{gsm8}k with a minimal (0.1$\%$) effect on accuracy.
    Our work demonstrates that the training objective and compute used by \glspl{mdlm} unlock many new possibilities for inference and post-training methods. Code is available \href{https://github.com/rajesh-lab/ReasoningDiffusionModels}{here}. \def\thefootnote{*}\footnotetext{Denotes equal authorship. Correspondence to \url{zfh2000@columbia.edu}, \url{rsinghal@nyu.edu}.}
\end{abstract}
